\documentclass[10pt,a4paper]{amsart}
\usepackage[margin=19mm]{geometry}
\usepackage{amsmath,amssymb,mathtools}
\usepackage[scr=rsfs,cal=euler]{mathalfa}
\usepackage{xfrac}
\usepackage[unicode,hidelinks,bookmarks=false]{hyperref}
\newcommand{\ie}{i.e.\ }
\newcommand{\cf}{cf.\ }
\newcommand{\resp}{respectively\ }
\newcommand{\Subsection}{\subsection}
\providecommand{\nobreakdash}{\nobreak\hbox{-}}
\setlength{\emergencystretch}{2em}
\multlinegap0pt
\usepackage{braket}
\usepackage{amssymb}
\usepackage{tcolorbox}
\usepackage{csquotes}
\usepackage{mathtools}
\usepackage{float}
\newtheorem{lemma}{Lemma}
\newtheorem{proposition}{Proposition}
\newcommand{\Hom}{\operatorname{Hom}}
\newcommand{\Ind}{\operatorname{Ind}}
\newcommand{\Res}{\operatorname{Res}}
\newcommand{\calC}{\mathcal{C}}
\newcommand{\calH}{\mathcal{H}}
\newcommand{\id}{\operatorname{id}}
\def\ket#1{| #1 \rangle}
\title{Hybrid Clifford Codes via Operator Algebra Quantum Error Correction and Projective Representation Theory}
\author{Jonas Eidesen, David W. Kribs, Andrew Nemec}
\date{Source: arXiv:2606.02531v1 (CC BY 4.0)}
\begin{document}
\maketitle

\noindent\emph{Source and licence.} Hybrid Clifford Codes via Operator Algebra Quantum Error Correction and Projective Representation Theory. Version arXiv:2606.02531v1 (CC BY 4.0), distributed by arXiv under the Creative Commons Attribution 4.0 International licence, http://creativecommons.org/licenses/by/4.0/, as recorded in the arXiv licence metadata for this exact version and shown on the abstract page https://arxiv.org/abs/2606.02531v1. Full licence notice: \texttt{licenses/arxiv-2606.02531-CC-BY-4.0.txt}.\par\smallskip

\noindent\textit{Editorial scope.} Counted unit: the article's proposition projection (source0.tex lines 771--776). The article's complete original proof of it is delivered with it (source0.tex lines 777--815, one \texttt{proof} environment of 39 lines). No mathematics is written, repaired or re-derived: the statement and the proof are the article's own lines, in the article's own notation, and no retained line is substituted or re-flowed.  Retained uncounted as the source-local context the proof needs: lines 224--226; lines 465--467; lines 745--754.  Nothing was omitted from any retained span, so the package carries no omission record.  No retained line contains a citation, so the wrapper carries no bibliography and no bibliography entry.  No retained line contains a figure environment, an image-inclusion command or drawing code, so no image is delivered and the package declares no asset dependency.  Editorial formatting only, in addition: an AMS \texttt{amsart} preamble that restates the article's own declarations and macros used by the retained lines, namely the environments \texttt{lemma}, \texttt{proposition} on separate counters, together with the standard packages those lines need.  The retained environments print in this document as Lemma 1, Proposition 1 in order of appearance, whereas the article prints the same items with the numbering of its own full document.  The complete native source of the article as distributed in the arXiv e-print for this exact version is bundled unchanged as \texttt{sources/201910/source0.tex}.
\begin{equation}\label{eqn:inner porduct of characters}
    \langle \chi_{\pi_1}, \chi_{\pi_2} \rangle = \dim \Hom_G(\pi_1,\pi_2).
\end{equation}

    \begin{equation}\label{eqn:clifford code representation}
        \Res^G_L\pi|_\calC = \gamma.
    \end{equation}

\begin{lemma}\label{lemma:intertwiners from class functions}
    Let $G$ be a finite group and $\pi$ be a $\sigma$-projective representation of $G$ on $\calH$. If $f$ is a $\sigma$-class function, then the linear map $T_f \colon \calH \to \calH$ defined by
    \begin{equation*}
        T_f = \sum_{x \in G}\overline{f(x)}\pi(x),
    \end{equation*}
    is an intertwiner. Furthermore, if $\pi$ is irreducible,
    \begin{equation*}
        T_f = \frac{|G|}{\dim\calH}\langle \chi_\pi, f \rangle \id_\calH.
    \end{equation*}
\end{lemma}

\begin{proposition}\label{lemma:code projection}
    Let $\calH$ be a finite dimensional Hilbert space and $(G,\pi)$ be a projective error model on $\calH$. Suppose that $\calC = \calC(G,\pi,L,\gamma)$ is a Clifford subspace code. Then the orthogonal projection onto $\calC$ is given by
    \begin{equation*}
        P_\calC = \frac{\dim\calC}{|L|}\sum_{x \in L}\overline{\chi_\gamma(x)}\pi(x).
    \end{equation*}
\end{proposition}

\begin{proof}
    To simplify notation let
    \begin{equation*}
        T = \frac{\dim\calC}{|L|}\sum_{x \in L}\overline{\chi_\gamma(x)}\pi(x)
    \end{equation*}
    and write $\calH = \calC \oplus \calC^\perp$. To show that $T = P_\calC$ it is enough to show that $T\ket{\psi} = \ket{\psi}$ for any $\ket{\psi} \in \calC$ and that $T\ket{\varphi} = 0$ for any $\ket{\varphi} \in \calC^\perp$. To show that the first holds, let $\ket{\psi} \in \calC$. Then,
    \begin{align*}
        T \ket{\psi}
        & = \frac{\dim\calC}{|L|}\sum_{x \in L}\overline{\chi_\gamma(x)}\pi(x)\ket{\psi} & \\
        & = \frac{\dim\calC}{|L|}\sum_{x \in L}\overline{\chi_\gamma(x)}\gamma(x)\ket{\psi} & \text{since } \pi(x)|_\calC = \gamma(x), \text{ see \eqref{eqn:clifford code representation}}, \\
        & = \langle\chi_\gamma, \chi_\gamma\rangle\ket{\psi} & \text{by Lemma~\ref{lemma:intertwiners from class functions}}, \\
        & = \ket{\psi} & \text{by Schur's Lemma and \eqref{eqn:inner porduct of characters}}.
    \end{align*}

    Now let $\ket{\varphi} \in \calC^\perp$ and let $V$ denote the smallest $L$-invariant subspace of $\calH$ containing $\ket{\varphi}$. By Maschke's Theorem we may assume that $\Res^G_L\pi|_V$ is an irreducible projective representation of $L$ on $V$. Let $\delta := \Res^G_L\pi|_V$, then by the same computation as above we see that
    \begin{equation*}
        T \ket{\varphi} = \frac{\dim\calC}{\dim V}\langle\chi_\delta, \chi_\gamma\rangle\ket{\varphi}.
    \end{equation*}
    By Schur's Lemma we have that
    \begin{equation*}
        \langle\chi_\delta, \chi_\gamma\rangle =
        \begin{cases}
            1, & \delta \simeq \gamma, \\
            0, & \delta \not\simeq \gamma.
        \end{cases}
    \end{equation*}
    By construction we have that $\delta \neq \gamma$, but they might still be isomorphic. To see that they cannot be isomorphic, notice that both are subrepresentations of $\Res^G_L\pi$. Combining Maschke's Theorem and Schur's Lemma we have that the number of subrepresentations of $\Res^G_L\pi$ that are isomorphic to $\gamma$ equals
    \begin{equation*}
        \dim \Hom_L(\gamma, \Res^G_L\pi).
    \end{equation*}
    We compute that
    \begin{align*}
        \dim \Hom_L(\gamma, \Res^G_L\pi)
        & = \dim \Hom_G(\Ind_L^G\gamma, \pi) & \text{by Frobenius reciprocity}, \\
        & = \dim \Hom_G(\pi, \pi) & \text{since $\calC$ is a Clifford subspace code}, \\
        & = 1 & \text{by Schur's Lemma}.
    \end{align*}
    Hence, $\gamma$ is the only subrepresentation of $\Res^G_L\pi$ that is isomorphic to $\gamma$, from which we conclude that $T\ket{\varphi} = 0$.
\end{proof}

\end{document}
